% ==========================================================================
%  NOTE: This is NOT the actual thesis. Working notes / draft material on the
%  feasibility of bringing the BioDT Honeybee Shiny application into the
%  ShinySwarm architecture. Scratch prose to review and selectively pull into
%  the real thesis (Future Work / Discussion). Verify all repo facts against
%  the live repositories before citing -- they were inspected on 2026-06-12.
% ==========================================================================

\documentclass[12pt,a4paper]{article}
\usepackage[utf8]{inputenc}
\usepackage[T1]{fontenc}
\usepackage{geometry}
\geometry{margin=2.5cm}
\usepackage{hyperref}
\usepackage{booktabs}
\usepackage{enumitem}
\usepackage{xcolor}

\title{\textbf{Feasibility: BioDT Honeybee in ShinySwarm}\\[0.5em]
\large Working notes on re-engineering the deferred use case}
\author{Iva Gunjaca\\
Master of Software Engineering, University of Amsterdam}
\date{\today}

\begin{document}
\maketitle

% --------------------------------------------------------------------------
\section{Why this section exists}
% --------------------------------------------------------------------------

The original proposal named a concrete scientific use case --- the
\textbf{BioDT Honeybee pDT / Pollinators UC} Shiny application --- as the
real-world subject onto which the ShinySwarm collaborative architecture would be
applied. That re-engineering was \emph{deferred to future work}, and a suite of
eight purpose-built Shiny apps was used for the benchmark instead, to gain full
control over experimental variables. This section assesses whether, and how, the
BioDT Honeybee application could be brought into ShinySwarm --- i.e.\ the
feasibility of completing the originally proposed use case.

\textbf{Bottom line up front.} It is feasible, and the application is a
\emph{better} architectural fit for ShinySwarm than the proposal assumed,
because the BioDT team has \emph{already} moved the compute to a microservice
(a containerised NetLogo/BEEHAVE job, optionally launched as a Kubernetes Job).
The remaining work is not ``decompose a monolith'' --- that is largely done ---
but ``add the collaboration and shared-state layer,'' which is exactly what
ShinySwarm contributes. The principal mismatch is framework-level: the BioDT app
is built on \textbf{Rhino} with \texttt{box} modules, whereas the ShinySwarm
apps are plain Shiny; this affects how state capture (\texttt{statesnap}) is
wired, not whether it can be done.

% --------------------------------------------------------------------------
\section{What the application actually is}
% --------------------------------------------------------------------------

The naming in the proposal is slightly misleading, so it is worth pinning down
the real artefacts. There is no single ``BioDT Honeybee Shiny app'' repository;
the use case is split across several BioDT GitHub repositories:

\begin{itemize}[nosep]
  \item \textbf{\texttt{BioDT/biodt-shiny}} --- the actual Shiny front end. A
    large, actively maintained (last push 2026-04), \textbf{Rhino}-framework
    application (\textasciitilde 630~kB of R) that hosts \emph{multiple} BioDT
    use cases as tabs/modules: honeybee, grassland, forest, disease outbreaks,
    crop wild relatives (cwr), cultural ecosystem services (ces), invasive alien
    species (ias), real-time bird monitoring (rtbm). The honeybee module is one
    pDT among many. License: ``Other'' (check before reuse).
  \item \textbf{\texttt{BioDT/uc-pollinators}} --- \emph{not} a Shiny app. The
    honeybee simulation pipeline: R scripts in three steps (prepare habitat
    JSONs $\rightarrow$ build BEEHAVE input $\rightarrow$ run BEEHAVE), packaged
    in a Docker image with NetLogo~6.3 + OpenJDK~17 + HyperQueue, designed to run
    on HPC (LUMI) or in the cloud. This is what the proposal called the
    ``Honeybee application,'' but it is the compute backend, not the UI. MIT.
  \item \textbf{\texttt{BioDT/uc-pollinators-beehave}} (NetLogo model) and
    \textbf{\texttt{BioDT/uc-pollinators-container}} (Docker build) --- git
    submodules of \texttt{uc-pollinators} holding the model and the container
    definition.
\end{itemize}

So the ``app'' the thesis would re-engineer is the \textbf{honeybee module of
\texttt{biodt-shiny}}, which orchestrates the \texttt{uc-pollinators} compute
container. The proposal's mention of components by \emph{Sandro Boelsz} refers to
this pipeline work.

% --------------------------------------------------------------------------
\section{Current architecture of the honeybee module}
% --------------------------------------------------------------------------

The honeybee module is already cleanly decomposed (Rhino \texttt{box} modules
under \texttt{app/view/honeybee/} and \texttt{app/logic/honeybee/}):

\begin{itemize}[nosep]
  \item \texttt{honeybee\_beekeeper.R} --- parent module wiring the sub-modules
    together and holding the shared reactive values (\texttt{map},
    \texttt{leaflet\_map}, \texttt{lookup\_table}, \texttt{experiment\_list}).
  \item \texttt{beekeeper\_map.R} --- Leaflet map; the user picks a location
    (\texttt{coordinates()}).
  \item \texttt{beekeeper\_param.R} --- simulation parameters.
  \item \texttt{beekeeper\_lookup.R} --- editable lookup table (habitat $\to$
    forage).
  \item \texttt{beekeeper\_runsimulation.R} --- writes the inputs (map crop,
    lookup, parameters, locations as CSV/GeoTIFF) into a per-run directory and
    \emph{launches the compute container}.
  \item \texttt{beekeeper\_plot.R} --- plots the result CSV.
  \item \texttt{honeybee/k8s.R} --- creates a Kubernetes \texttt{batch/v1} Job
    that runs \texttt{ghcr.io/biodt/beehave:0.3.13}, watches it to completion,
    then deletes it.
\end{itemize}

Crucially, the run step already supports \textbf{two executors}, selected by
\texttt{config\$get("executor")}:
\begin{enumerate}[nosep]
  \item \texttt{"docker"} --- shells out with \texttt{system()} to
    \texttt{docker run ... ghcr.io/biodt/beehave:0.3.13}, bind-mounting a run
    directory.
  \item \texttt{"k8s"} --- POSTs a Job spec to the in-cluster Kubernetes API
    (token from the mounted service account), mounts per-user PVCs keyed by
    \texttt{SHINYPROXY\_ID}, waits via a watch stream, deletes the Job.
\end{enumerate}

Deployment today is \textbf{ShinyProxy on Kubernetes}: one container per user
session (hence \texttt{SHINYPROXY\_ID} and per-user PVC names), with the heavy
simulation offloaded to ephemeral Job pods. State is held in per-session reactive
values and in a per-session ``shared'' directory on disk; there is \emph{no}
cross-user sharing and \emph{no} real-time collaboration.

\textbf{Implication.} BioDT has independently arrived at part of ShinySwarm's
thesis: containerise the scientific compute and orchestrate it on Kubernetes.
What it lacks is precisely ShinySwarm's contribution --- multi-user, real-time
\emph{collaboration} and portable, reproducible \emph{shared state}.

% --------------------------------------------------------------------------
\section{Feasibility of bringing it into ShinySwarm}
% --------------------------------------------------------------------------

Mapping the honeybee module onto ShinySwarm's two architectures (REST API with
Plumber + Redis; or Kafka event-driven) is feasible. The work decomposes into
four parts, in rough order of effort.

\subsection{1. Compute service --- already done (low effort)}

ShinySwarm's premise is that R computation lives behind a service. For honeybee
this already exists: \texttt{ghcr.io/biodt/beehave} is a self-contained compute
container with a CSV/GeoTIFF-in, CSV-out contract. It can be driven unchanged.
The only adaptation is \emph{how it is invoked}:
\begin{itemize}[nosep]
  \item In the REST variant, wrap the existing Docker/k8s launch behind a Plumber
    endpoint (\texttt{POST /simulate} with the input bundle, returns a run id;
    \texttt{GET /result/\{id\}}), instead of calling \texttt{system()} from the
    UI session. The launch logic in \texttt{beekeeper\_runsimulation.R} and
    \texttt{k8s.R} moves almost verbatim into the service.
  \item In the Kafka variant, the UI produces a ``simulate'' message to a topic;
    a consumer launches the BEEHAVE Job and produces a ``result-ready'' message
    with the output location. This matches BEEHAVE's 2--4~minute, long-running,
    asynchronous nature far better than a synchronous request, and is a strong
    argument for the event-driven architecture on this use case.
\end{itemize}

\subsection{2. Shared state and collaboration --- the real work (medium)}

This is what ShinySwarm adds and the BioDT app lacks. The honeybee module's
collaborative state is small and well-defined, which helps:
\begin{itemize}[nosep]
  \item \textbf{Inputs to share:} selected coordinates, parameters, the editable
    lookup table, and the chronological \texttt{experiment\_list} of runs.
  \item \textbf{Outputs to share:} the result CSV(s) and the rendered plot/map.
\end{itemize}
Under ShinySwarm these become shared state synchronised through Redis (REST) or
Kafka topics (event-driven), so that multiple beekeepers/researchers see the same
map selection, parameters, and run history in real time --- the live ``shared
state management'' the proposal promised.

\subsection{3. Reproducible state capture with \texttt{statesnap} (low--medium)}

The honeybee module is an almost ideal showcase for \texttt{statesnap}, because
its result is the output of a \textbf{non-deterministic BEEHAVE simulation}.
Re-running from the same inputs need not reproduce the same population trajectory,
so input-only sharing (Shiny bookmarking) cannot reproduce a colleague's exact
result --- whereas \texttt{statesnap}, which captures the \emph{computed} result,
can. Concretely:
\begin{itemize}[nosep]
  \item Register the module's reactives by name (\texttt{coordinates},
    \texttt{parameters}, \texttt{lookup}, \texttt{experiment\_list}).
  \item Embed the input map crop and the result CSV with \texttt{state\_file()};
    the bundle is self-contained.
  \item A captured checkpoint becomes a shareable, reproducible ``saved
    exploration'' --- the cross-user restore the benchmark's RQ5 already exercises
    on the eight custom apps, now demonstrated on a real scientific application.
\end{itemize}
\textbf{Caveat (framework):} \texttt{statesnap} currently reads
\texttt{names(input)} at the top level and pushes values through the parent
session; it does not yet handle Shiny \emph{module namespaces}. The honeybee UI is
deeply modular (Rhino \texttt{box} + nested \texttt{moduleServer}). So either
(a)~capture/restore is invoked \emph{inside} each module's own server and the
fragments merged, or (b)~\texttt{statesnap} gains namespace-aware capture (already
listed as further work). This is the single most important technical dependency
and should be stated as such.

\subsection{4. Framework reconciliation --- Rhino vs plain Shiny (medium)}

The ShinySwarm benchmark apps are plain Shiny; \texttt{biodt-shiny} is Rhino with
\texttt{box} imports, \texttt{config} profiles, an \texttt{i18n} translation
layer, SCSS/JS build via webpack, and Cypress tests. None of this blocks
integration --- Rhino apps are still Shiny --- but it means the integration cannot
be a copy-paste of the benchmark harness. Options:
\begin{itemize}[nosep]
  \item \textbf{Extract} the honeybee module into a standalone Shiny app and wrap
    it in ShinySwarm (cleanest for a controlled experiment; loses the multi-pDT
    context).
  \item \textbf{Integrate in place} --- add the ShinySwarm REST/Kafka client and
    shared-state layer as Rhino \texttt{logic} modules inside \texttt{biodt-shiny}
    (most faithful to the real system; more coupling to BioDT's build).
\end{itemize}

\subsection{Other integration considerations}

\begin{itemize}[nosep]
  \item \textbf{Per-session vs shared model.} BioDT uses ShinyProxy
    one-container-per-user with per-user PVCs. ShinySwarm's collaborative model
    deliberately shares state across users; the deployment topology (shared
    container vs container-per-user) is itself one of the architectural variables
    the proposal set out to study, so honeybee is a natural place to revisit it.
  \item \textbf{Auth.} \texttt{biodt-shiny} runs behind ShinyProxy; ShinySwarm
    uses JWT (Spring Security). A real deployment would reconcile these (and,
    longer term, a shared Keycloak/OIDC layer as discussed in the NaaVRE notes).
  \item \textbf{Data dependencies.} The module needs habitat map tiles, a lookup
    table, and (for weather) Copernicus CDS / HDA credentials. These must be
    provisioned to the compute service rather than the UI.
  \item \textbf{Licensing.} \texttt{uc-pollinators} is MIT; \texttt{biodt-shiny}
    is ``Other'' --- confirm terms before redistributing modified code.
\end{itemize}

% --------------------------------------------------------------------------
\section{Effort summary and risk}
% --------------------------------------------------------------------------

\begin{table}[h]
\centering
\small
\begin{tabular}{@{}lll@{}}
\toprule
\textbf{Work item} & \textbf{Effort} & \textbf{Risk / dependency} \\
\midrule
Drive existing BEEHAVE container via service & Low & Contract already CSV/CSV \\
REST (Plumber+Redis) wrapper & Low--Med & Reuse existing launch code \\
Kafka async wrapper & Medium & Best fit for 2--4 min jobs \\
Shared state + real-time sync & Medium & Core ShinySwarm contribution \\
\texttt{statesnap} reproducible checkpoints & Low--Med & \emph{Needs module-namespace support} \\
Rhino vs plain-Shiny reconciliation & Medium & Build/i18n/test coupling \\
Auth (ShinyProxy $\to$ JWT/OIDC) & Medium & Cross-system identity \\
\bottomrule
\end{tabular}
\caption{Feasibility breakdown for the BioDT Honeybee use case in ShinySwarm.}
\end{table}

\textbf{Conclusion.} Re-engineering the BioDT Honeybee use case into ShinySwarm
is feasible and would strengthen the thesis by validating the architecture on a
real, non-deterministic scientific application rather than synthetic apps. The
compute is already a container, so the proposal's hardest assumed task is
effectively pre-solved; the substantive work is the collaboration and shared-state
layer that ShinySwarm exists to provide. The two genuine prerequisites are
(i)~\texttt{statesnap} module-namespace support, needed for reproducible state
capture in the deeply modular Rhino UI, and (ii)~a decision between extracting the
honeybee module versus integrating ShinySwarm in place. Neither is a blocker;
both are scoped, known pieces of work --- which is precisely why the
re-engineering was a reasonable item to defer to future work rather than abandon.

% --------------------------------------------------------------------------
\section{Source provenance}
% --------------------------------------------------------------------------

Repositories inspected on 2026-06-12 (GitHub API + raw file reads):
\begin{itemize}[nosep]
  \item \url{https://github.com/BioDT/biodt-shiny} (honeybee module, config.yml,
    docker-compose.yml, k8s.R, runsimulation module).
  \item \url{https://github.com/BioDT/uc-pollinators} (compute pipeline,
    Dockerfile, Makefile, submodules).
  \item Proposal reference [5]: BioDT, ``uc-pollinators: Honeybee pDT/pollinators
    UC main repository,'' 2024.
\end{itemize}

\end{document}
